\documentclass[12pt,a4paper]{article}
\usepackage{fontspec}\setmainfont{Latin Modern Roman}\setsansfont{DejaVu Sans}
\usepackage{amsmath,amssymb,graphicx,geometry,enumitem,needspace,longtable,booktabs,fancyhdr,hyperref}
\geometry{left=28mm,right=23mm,top=24mm,bottom=24mm}
\hypersetup{unicode=true,colorlinks=true,linkcolor=black,urlcolor=blue,pdftitle={Giải tích §7.3 — Phép thế lượng giác}}
\setlength{\parindent}{1.2em}\setlength{\parskip}{4pt}\setlength{\headheight}{16pt}
\pagestyle{fancy}\fancyhf{}\fancyhead[L]{\small Chương 7. Các phương pháp tính tích phân}\fancyhead[R]{\small §7.3}\fancyfoot[C]{\thepage}
\newcommand{\fig}[3][.68]{\begin{center}\begin{minipage}{\linewidth}\centering\includegraphics[width=#1\linewidth]{assets/#2.png}\\{\small #3}\end{minipage}\end{center}}
\newcommand{\ex}[1]{\par\medskip\noindent\textbf{Ví dụ #1.}\ }
\newcommand{\sol}{\par\nopagebreak[4]\noindent\textit{Lời giải.}\ }
\newcommand{\bt}[1]{\par\smallskip\Needspace{5\baselineskip}\noindent\textbf{#1.}\ }
\newcommand{\group}[1]{\par\Needspace{11\baselineskip}\medskip\noindent\textbf{#1}\par\nopagebreak[4]}
\newcommand{\gint}[1]{\left[\!\left[#1\right]\!\right]}
\allowdisplaybreaks[2]
\DeclareMathOperator{\sech}{sech}
\DeclareMathOperator{\csch}{csch}
\DeclareMathOperator{\arccot}{arccot}

\newcommand{\prj}[1]{\bt{#1}}

\newcommand{\kn}[1]{\bt{#1}}
\newcommand{\ds}[1]{\bt{#1}}
\setlength{\emergencystretch}{1em}
\begin{document}
\begin{center}{\LARGE\bfseries 7.3 — Phép thế lượng giác}\end{center}
Khi tìm diện tích của một hình tròn hoặc một elip, ta gặp một tích phân dạng $\int\sqrt{a^2-x^2}\,dx$, trong đó $a>0$. Nếu tích phân là $\int x\sqrt{a^2-x^2}\,dx$ thì phép đổi biến $u=a^2-x^2$ có hiệu quả, nhưng tích phân $\int\sqrt{a^2-x^2}\,dx$ như đang xét thì khó hơn. Nếu đổi biến từ $x$ sang $\theta$ bằng phép thế $x=a\sin\theta$, đồng nhất thức $1-\sin^2\theta=\cos^2\theta$ cho phép loại bỏ dấu căn vì
\begin{align*}
\sqrt{a^2-x^2}&=\sqrt{a^2-a^2\sin^2\theta}=\sqrt{a^2(1-\sin^2\theta)}\\
&=\sqrt{a^2\cos^2\theta}=a|\cos\theta|.
\end{align*}
Chú ý sự khác nhau giữa phép đổi biến $u=a^2-x^2$ (biến mới là một hàm của biến cũ) và phép thế $x=a\sin\theta$ (biến cũ là một hàm của biến mới).

Nói chung, ta có thể thực hiện phép thế dạng $x=g(t)$ bằng cách dùng quy tắc đổi biến theo chiều ngược lại. Để việc tính toán đơn giản hơn, ta giả sử $g$ có hàm ngược, tức là $g$ đơn ánh. Khi đó, nếu thay $u$ bằng $x$ và $x$ bằng $t$ trong quy tắc đổi biến (Phương trình 5.5.4), ta được
\[\int f(x)\,dx=\int f(g(t))g'(t)\,dt.\]
Kiểu đổi biến này được gọi là \textit{phép thế ngược}.

Ta có thể thực hiện phép thế ngược $x=a\sin\theta$ miễn là nó xác định một hàm đơn ánh. Có thể bảo đảm điều này bằng cách giới hạn $\theta$ trong đoạn $[-\pi/2,\pi/2]$.

Bảng dưới đây liệt kê những phép thế lượng giác có hiệu quả đối với các biểu thức chứa căn đã cho, nhờ các đồng nhất thức lượng giác tương ứng. Trong mỗi trường hợp, giới hạn đặt lên $\theta$ nhằm bảo đảm hàm xác định phép thế là đơn ánh. (Đây cũng là những khoảng đã dùng trong Mục 1.5 khi định nghĩa các hàm ngược.)

\begin{center}\textbf{Bảng các phép thế lượng giác}\end{center}
\begin{longtable}{p{.17\linewidth}p{.47\linewidth}p{.25\linewidth}}
\toprule
Biểu thức & Phép thế & Đồng nhất thức\\\midrule
$\sqrt{a^2-x^2}$ & $x=a\sin\theta$, $-\dfrac\pi2\le\theta\le\dfrac\pi2$ & $1-\sin^2\theta=\cos^2\theta$\\[8pt]
$\sqrt{a^2+x^2}$ & $x=a\tan\theta$, $-\dfrac\pi2<\theta<\dfrac\pi2$ & $1+\tan^2\theta=\sec^2\theta$\\[8pt]
$\sqrt{x^2-a^2}$ & $x=a\sec\theta$, $0\le\theta<\dfrac\pi2$ hoặc $\pi\le\theta<\dfrac{3\pi}2$ & $\sec^2\theta-1=\tan^2\theta$\\
\bottomrule
\end{longtable}

\ex{1} Tính $\displaystyle\int\frac{\sqrt{9-x^2}}{x^2}\,dx$.
\sol Đặt $x=3\sin\theta$, trong đó $-\pi/2\le\theta\le\pi/2$. Khi đó $dx=3\cos\theta\,d\theta$ và
\[\sqrt{9-x^2}=\sqrt{9-9\sin^2\theta}=\sqrt{9\cos^2\theta}=3|\cos\theta|=3\cos\theta.\]
(Chú ý rằng $\cos\theta\ge0$ vì $-\pi/2\le\theta\le\pi/2$.) Vì vậy, quy tắc thế ngược cho
\begin{align*}
\int\frac{\sqrt{9-x^2}}{x^2}\,dx&=\int\frac{3\cos\theta}{9\sin^2\theta}\,3\cos\theta\,d\theta\\
&=\int\frac{\cos^2\theta}{\sin^2\theta}\,d\theta=\int\cot^2\theta\,d\theta\\
&=\int(\csc^2\theta-1)\,d\theta\\
&=-\cot\theta-\theta+C.
\end{align*}
Vì đây là tích phân bất định, ta phải trở về biến ban đầu $x$. Có thể làm điều này bằng cách dùng các đồng nhất thức lượng giác để biểu diễn $\cot\theta$ theo $\sin\theta=x/3$, hoặc vẽ hình như Hình 1, trong đó $\theta$ được hiểu là một góc của tam giác vuông. Vì $\sin\theta=x/3$, ta ghi độ dài cạnh đối và cạnh huyền lần lượt là $x$ và $3$. Khi đó định lý Pythagore cho độ dài cạnh kề là $\sqrt{9-x^2}$, nên có thể đọc ngay giá trị của $\cot\theta$ từ hình:
\[\cot\theta=\frac{\sqrt{9-x^2}}x.\]
\fig[.34]{figure01}{Hình 1. $\sin\theta=x/3$.}
(Mặc dù $\theta>0$ trong hình, biểu thức này của $\cot\theta$ vẫn đúng khi $\theta<0$.) Vì $\sin\theta=x/3$, ta có $\theta=\sin^{-1}(x/3)$, do đó
\begin{align*}
\int\frac{\sqrt{9-x^2}}{x^2}\,dx&=-\cot\theta-\theta+C\\
&=-\frac{\sqrt{9-x^2}}x-\sin^{-1}\!\left(\frac x3\right)+C.
\end{align*}

\ex{2} Tìm diện tích miền bên trong elip
\[\frac{x^2}{a^2}+\frac{y^2}{b^2}=1.\]
\sol Giải phương trình elip theo $y$, ta được
\[\frac{y^2}{b^2}=1-\frac{x^2}{a^2}=\frac{a^2-x^2}{a^2},\qquad\text{hay}\qquad y=\pm\frac ba\sqrt{a^2-x^2}.\]
Vì elip đối xứng qua cả hai trục, tổng diện tích $A$ bằng bốn lần diện tích trong góc phần tư thứ nhất (xem Hình 2). Phần elip trong góc phần tư thứ nhất được cho bởi hàm
\[y=\frac ba\sqrt{a^2-x^2},\qquad0\le x\le a,\]
và vì vậy
\[\frac14A=\int_0^a\frac ba\sqrt{a^2-x^2}\,dx.\]
\fig[.34]{figure02}{Hình 2. $x^2/a^2+y^2/b^2=1$.}
Để tính tích phân này, ta đặt $x=a\sin\theta$. Khi đó $dx=a\cos\theta\,d\theta$. Để đổi cận tích phân, chú ý rằng khi $x=0$ thì $\sin\theta=0$, nên $\theta=0$; khi $x=a$ thì $\sin\theta=1$, nên $\theta=\pi/2$. Ngoài ra,
\[\sqrt{a^2-x^2}=\sqrt{a^2-a^2\sin^2\theta}=\sqrt{a^2\cos^2\theta}=a|\cos\theta|=a\cos\theta,\]
vì $0\le\theta\le\pi/2$. Do đó
\begin{align*}
A&=4\frac ba\int_0^a\sqrt{a^2-x^2}\,dx=4\frac ba\int_0^{\pi/2}a\cos\theta\cdot a\cos\theta\,d\theta\\
&=4ab\int_0^{\pi/2}\cos^2\theta\,d\theta=4ab\int_0^{\pi/2}\frac12(1+\cos2\theta)\,d\theta\\
&=2ab\left[\theta+\frac12\sin2\theta\right]_0^{\pi/2}=2ab\left(\frac\pi2+0-0\right)=\pi ab.
\end{align*}
Ta đã chứng minh rằng diện tích của một elip có các bán trục $a$ và $b$ là $\pi ab$. Đặc biệt, lấy $a=b=r$, ta đã chứng minh công thức quen thuộc: diện tích hình tròn bán kính $r$ là $\pi r^2$.

\textbf{Chú ý.} Vì tích phân trong Ví dụ 2 là một tích phân xác định, ta đã đổi cận tích phân và không cần chuyển trở lại biến ban đầu $x$.

\ex{3} Tìm $\displaystyle\int\frac{dx}{x^2\sqrt{x^2+4}}$.
\sol Đặt $x=2\tan\theta$, $-\pi/2<\theta<\pi/2$. Khi đó $dx=2\sec^2\theta\,d\theta$ và
\[\sqrt{x^2+4}=\sqrt{4(\tan^2\theta+1)}=\sqrt{4\sec^2\theta}=2|\sec\theta|=2\sec\theta.\]
Vì vậy, ta có
\[\int\frac{dx}{x^2\sqrt{x^2+4}}=\int\frac{2\sec^2\theta\,d\theta}{4\tan^2\theta\cdot2\sec\theta}=\frac14\int\frac{\sec\theta}{\tan^2\theta}\,d\theta.\]
Để tính tích phân lượng giác này, ta viết mọi thứ theo $\sin\theta$ và $\cos\theta$:
\[\frac{\sec\theta}{\tan^2\theta}=\frac1{\cos\theta}\cdot\frac{\cos^2\theta}{\sin^2\theta}=\frac{\cos\theta}{\sin^2\theta}.\]
Do đó, thực hiện phép đổi biến $u=\sin\theta$, ta có
\begin{align*}
\int\frac{dx}{x^2\sqrt{x^2+4}}&=\frac14\int\frac{\cos\theta}{\sin^2\theta}\,d\theta=\frac14\int\frac{du}{u^2}\\
&=\frac14\left(-\frac1u\right)+C=-\frac1{4\sin\theta}+C\\
&=-\frac{\csc\theta}{4}+C.
\end{align*}
Ta dùng Hình 3 để xác định rằng $\csc\theta=\sqrt{x^2+4}/x$, nên
\[\int\frac{dx}{x^2\sqrt{x^2+4}}=-\frac{\sqrt{x^2+4}}{4x}+C.\]
\fig[.34]{figure03}{Hình 3. $\tan\theta=x/2$.}

\ex{4} Tìm $\displaystyle\int\frac{x}{\sqrt{x^2+4}}\,dx$.
\sol Ở đây có thể dùng phép thế lượng giác $x=2\tan\theta$ (như trong Ví dụ 3). Nhưng phép đổi biến trực tiếp $u=x^2+4$ đơn giản hơn, vì khi đó $du=2x\,dx$ và
\[\int\frac{x}{\sqrt{x^2+4}}\,dx=\frac12\int\frac{du}{\sqrt u}=\sqrt u+C=\sqrt{x^2+4}+C.\]
\textbf{Chú ý.} Ví dụ 4 minh họa rằng ngay cả khi có thể thực hiện phép thế lượng giác, phương pháp này cũng chưa chắc cho lời giải dễ nhất. Trước hết, bạn nên tìm một phương pháp đơn giản hơn.

\ex{5} Tính $\displaystyle\int\frac{dx}{\sqrt{x^2-a^2}}$, trong đó $a>0$.
\sol \textbf{Cách 1.} Đặt $x=a\sec\theta$, trong đó $0<\theta<\pi/2$ hoặc $\pi<\theta<3\pi/2$. Khi đó $dx=a\sec\theta\tan\theta\,d\theta$ và
\[\sqrt{x^2-a^2}=\sqrt{a^2(\sec^2\theta-1)}=\sqrt{a^2\tan^2\theta}=a|\tan\theta|=a\tan\theta.\]
Do đó
\begin{align*}
\int\frac{dx}{\sqrt{x^2-a^2}}&=\int\frac{a\sec\theta\tan\theta}{a\tan\theta}\,d\theta\\
&=\int\sec\theta\,d\theta=\ln|\sec\theta+\tan\theta|+C.
\end{align*}
Tam giác trong Hình 4 cho $\tan\theta=\sqrt{x^2-a^2}/a$, nên ta có
\begin{align*}
\int\frac{dx}{\sqrt{x^2-a^2}}&=\ln\left|\frac xa+\frac{\sqrt{x^2-a^2}}a\right|+C\\
&=\ln\left|x+\sqrt{x^2-a^2}\right|-\ln a+C.
\end{align*}
\fig[.32]{figure04}{Hình 4. $\sec\theta=x/a$.}
Viết $C_1=C-\ln a$, ta được
\begin{equation}\int\frac{dx}{\sqrt{x^2-a^2}}=\ln\left|x+\sqrt{x^2-a^2}\right|+C_1.\tag{1}\end{equation}
\sol \textbf{Cách 2.} Với $x>0$, cũng có thể dùng phép thế hyperbolic $x=a\cosh t$. Sử dụng đồng nhất thức $\cosh^2y-\sinh^2y=1$, ta có
\[\sqrt{x^2-a^2}=\sqrt{a^2(\cosh^2t-1)}=\sqrt{a^2\sinh^2t}=a\sinh t.\]
Vì $dx=a\sinh t\,dt$, ta được
\[\int\frac{dx}{\sqrt{x^2-a^2}}=\int\frac{a\sinh t\,dt}{a\sinh t}=\int dt=t+C.\]
Vì $\cosh t=x/a$, ta có $t=\cosh^{-1}(x/a)$ và
\begin{equation}\int\frac{dx}{\sqrt{x^2-a^2}}=\cosh^{-1}\!\left(\frac xa\right)+C.\tag{2}\end{equation}
Mặc dù Công thức 1 và 2 trông rất khác nhau, chúng thực ra tương đương theo Công thức 3.11.4.

\textbf{Chú ý.} Như Ví dụ 5 minh họa, có thể dùng phép thế hyperbolic thay cho phép thế lượng giác, và đôi khi chúng dẫn đến kết quả đơn giản hơn. Tuy nhiên, ta thường dùng phép thế lượng giác vì các đồng nhất thức lượng giác quen thuộc hơn các đồng nhất thức hyperbolic.

\ex{6} Tìm $\displaystyle\int_0^{3\sqrt3/2}\frac{x^3}{(4x^2+9)^{3/2}}\,dx$.
\sol Trước hết, chú ý rằng $(4x^2+9)^{3/2}=(\sqrt{4x^2+9})^3$, nên phép thế lượng giác là thích hợp. Mặc dù $\sqrt{4x^2+9}$ chưa hoàn toàn là một trong các biểu thức trong bảng phép thế lượng giác, nó trở thành một biểu thức như vậy nếu ta thực hiện phép đổi biến sơ bộ $u=2x$, cho $\sqrt{u^2+9}$. Sau đó đặt $u=3\tan\theta$, hay tương đương $x=\frac32\tan\theta$, suy ra $dx=\frac32\sec^2\theta\,d\theta$ và
\[\sqrt{4x^2+9}=\sqrt{9\tan^2\theta+9}=3\sec\theta.\]
Khi $x=0$ thì $\tan\theta=0$, nên $\theta=0$; khi $x=3\sqrt3/2$ thì $\tan\theta=\sqrt3$, nên $\theta=\pi/3$.
\begin{align*}
\int_0^{3\sqrt3/2}\frac{x^3}{(4x^2+9)^{3/2}}\,dx
&=\int_0^{\pi/3}\frac{\frac{27}{8}\tan^3\theta}{27\sec^3\theta}\,\frac32\sec^2\theta\,d\theta\\
&=\frac3{16}\int_0^{\pi/3}\frac{\tan^3\theta}{\sec\theta}\,d\theta
=\frac3{16}\int_0^{\pi/3}\frac{\sin^3\theta}{\cos^2\theta}\,d\theta\\
&=\frac3{16}\int_0^{\pi/3}\frac{1-\cos^2\theta}{\cos^2\theta}\sin\theta\,d\theta.
\end{align*}
Bây giờ đặt $u=\cos\theta$, nên $du=-\sin\theta\,d\theta$. Khi $\theta=0$ thì $u=1$; khi $\theta=\pi/3$ thì $u=\frac12$. Do đó
\begin{align*}
\int_0^{3\sqrt3/2}\frac{x^3}{(4x^2+9)^{3/2}}\,dx
&=-\frac3{16}\int_1^{1/2}\frac{1-u^2}{u^2}\,du\\
&=\frac3{16}\int_1^{1/2}(1-u^{-2})\,du
=\frac3{16}\left[u+\frac1u\right]_1^{1/2}\\
&=\frac3{16}\left[\left(\frac12+2\right)-(1+1)\right]=\frac3{32}.
\end{align*}
\textit{Ghi chú bên lề.} Như Ví dụ 6 cho thấy, phép thế lượng giác đôi khi là một lựa chọn tốt khi trong tích phân xuất hiện $(x^2+a^2)^{n/2}$, với $n$ là một số nguyên bất kỳ. Điều tương tự cũng đúng khi xuất hiện $(a^2-x^2)^{n/2}$ hoặc $(x^2-a^2)^{n/2}$.

\ex{7} Tính $\displaystyle\int\frac{x}{\sqrt{3-2x-x^2}}\,dx$.
\sol Ta có thể biến đổi hàm dưới dấu tích phân thành một hàm thích hợp với phép thế lượng giác bằng cách trước hết hoàn thành bình phương biểu thức dưới dấu căn:
\begin{align*}
3-2x-x^2&=3-(x^2+2x)=3+1-(x^2+2x+1)\\
&=4-(x+1)^2.
\end{align*}
Điều này gợi ý phép đổi biến $u=x+1$. Khi đó $du=dx$ và $x=u-1$, nên
\[\int\frac{x}{\sqrt{3-2x-x^2}}\,dx=\int\frac{u-1}{\sqrt{4-u^2}}\,du.\]
Bây giờ đặt $u=2\sin\theta$, suy ra $du=2\cos\theta\,d\theta$ và $\sqrt{4-u^2}=2\cos\theta$, do đó
\begin{align*}
\int\frac{x}{\sqrt{3-2x-x^2}}\,dx&=\int\frac{2\sin\theta-1}{2\cos\theta}\,2\cos\theta\,d\theta\\
&=\int(2\sin\theta-1)\,d\theta\\
&=-2\cos\theta-\theta+C\\
&=-\sqrt{4-u^2}-\sin^{-1}\!\left(\frac u2\right)+C\\
&=-\sqrt{3-2x-x^2}-\sin^{-1}\!\left(\frac{x+1}{2}\right)+C.
\end{align*}
Hình 5 cho thấy các đồ thị của hàm dưới dấu tích phân trong Ví dụ 7 và tích phân bất định của nó (với $C=0$). Đồ thị nào ứng với hàm nào?
\fig[.32]{figure05}{Hình 5.}

\clearpage
\section*{7.3 — Bài tập}
\group{1–4} (a) Xác định một phép thế lượng giác thích hợp. (b) Áp dụng phép thế để biến đổi tích phân thành một tích phân lượng giác. Không tính tích phân.
\bt{1} $\displaystyle\int\frac{x^3}{\sqrt{1+x^2}}\,dx$.
\bt{2} $\displaystyle\int\frac{x^3}{\sqrt{9-x^2}}\,dx$.
\bt{3} $\displaystyle\int\frac{x^2}{\sqrt{x^2-2}}\,dx$.
\bt{4} $\displaystyle\int\frac{x^3}{(9-4x^2)^{3/2}}\,dx$.

\group{5–8} Tính tích phân bằng phép thế lượng giác đã cho. Phác họa và ghi các kích thước, góc của tam giác vuông tương ứng.
\bt{5} $\displaystyle\int\frac{x^3}{\sqrt{1-x^2}}\,dx$; $x=\sin\theta$.
\bt{6} $\displaystyle\int\frac{x^3}{\sqrt{9+x^2}}\,dx$; $x=3\tan\theta$.
\bt{7} $\displaystyle\int\frac{\sqrt{4x^2-25}}x\,dx$; $x=\frac52\sec\theta$.
\bt{8} $\displaystyle\int\frac{\sqrt{2-x^2}}{x^2}\,dx$; $x=\sqrt2\sin\theta$.

\group{9–36} Tính tích phân.
\bt{9} $\displaystyle\int x^3\sqrt{16+x^2}\,dx$.
\bt{10} $\displaystyle\int\frac{x^2}{\sqrt{9-x^2}}\,dx$.
\bt{11} $\displaystyle\int\frac{\sqrt{x^2-1}}{x^4}\,dx$.
\bt{12} $\displaystyle\int_0^3\frac{x}{\sqrt{36-x^2}}\,dx$.
\bt{13} $\displaystyle\int_0^a\frac{dx}{(a^2+x^2)^{3/2}}$, $a>0$.
\bt{14} $\displaystyle\int\frac{dt}{t^2\sqrt{t^2-16}}$.
\bt{15} $\displaystyle\int_2^3\frac{dx}{(x^2-1)^{3/2}}$.
\bt{16} $\displaystyle\int_0^{2/3}\sqrt{4-9x^2}\,dx$.
\bt{17} $\displaystyle\int_0^{1/2}x\sqrt{1-4x^2}\,dx$.
\bt{18} $\displaystyle\int_0^2\frac{dt}{\sqrt{4+t^2}}$.
\bt{19} $\displaystyle\int\frac{\sqrt{x^2-9}}{x^3}\,dx$.
\bt{20} $\displaystyle\int_0^1\frac{dx}{(x^2+1)^2}$.
\bt{21} $\displaystyle\int_0^a x^2\sqrt{a^2-x^2}\,dx$.
\bt{22} $\displaystyle\int_{1/4}^{\sqrt3/4}\sqrt{1-4x^2}\,dx$.
\bt{23} $\displaystyle\int\frac{x}{\sqrt{x^2-7}}\,dx$.
\bt{24} $\displaystyle\int\frac{x}{\sqrt{1+x^2}}\,dx$.
\bt{25} $\displaystyle\int\frac{\sqrt{1+x^2}}x\,dx$.
\bt{26} $\displaystyle\int_0^{0.3}\frac{x}{(9-25x^2)^{3/2}}\,dx$.
\bt{27} $\displaystyle\int_0^{0.6}\frac{x^2}{\sqrt{9-25x^2}}\,dx$.
\bt{28} $\displaystyle\int_0^1\sqrt{x^2+1}\,dx$.
\bt{29} $\displaystyle\int\frac{dx}{\sqrt{x^2+2x+5}}$.
\bt{30} $\displaystyle\int_0^1\sqrt{x-x^2}\,dx$.
\bt{31} $\displaystyle\int x^2\sqrt{3+2x-x^2}\,dx$.
\bt{32} $\displaystyle\int\frac{x^2}{(3+4x-4x^2)^{3/2}}\,dx$.
\bt{33} $\displaystyle\int\sqrt{x^2+2x}\,dx$.
\bt{34} $\displaystyle\int\frac{x^2+1}{(x^2-2x+2)^2}\,dx$.
\bt{35} $\displaystyle\int x\sqrt{1-x^4}\,dx$.
\bt{36} $\displaystyle\int_0^{\pi/2}\frac{\cos t}{\sqrt{1+\sin^2t}}\,dt$.

\bt{37} (a) Dùng phép thế lượng giác để chứng minh rằng
\[\int\frac{dx}{\sqrt{x^2+a^2}}=\ln\left(x+\sqrt{x^2+a^2}\right)+C.\]
(b) Dùng phép thế hyperbolic $x=a\sinh t$ để chứng minh rằng
\[\int\frac{dx}{\sqrt{x^2+a^2}}=\sinh^{-1}\!\left(\frac xa\right)+C.\]
Hai công thức này liên hệ với nhau qua Công thức 3.11.3.

\bt{38} Tính
\[\int\frac{x^2}{(x^2+a^2)^{3/2}}\,dx\]
(a) bằng phép thế lượng giác;

(b) bằng phép thế hyperbolic $x=a\sinh t$.

\bt{39} Tìm giá trị trung bình của $f(x)=\sqrt{x^2-1}/x$, $1\le x\le7$.

\bt{40} Tìm diện tích miền giới hạn bởi hyperbol $9x^2-4y^2=36$ và đường thẳng $x=3$.

\par\noindent\begin{minipage}[t]{\linewidth}
\bt{41} Chứng minh công thức $A=\frac12r^2\theta$ cho diện tích hình quạt của một đường tròn có bán kính $r$ và góc ở tâm $\theta$. [Gợi ý: Giả sử $0<\theta<\pi/2$ và đặt tâm đường tròn tại gốc tọa độ, nên đường tròn có phương trình $x^2+y^2=r^2$. Khi đó $A$ bằng tổng diện tích tam giác $POQ$ và diện tích miền $PQR$ trong hình.]
\fig[.32]{exercise41}{Hình cho Bài tập 41.}

\end{minipage}\par
\bt{42} \textsf{[Đồ thị]} Tính tích phân
\[\int\frac{dx}{x^4\sqrt{x^2-2}}.\]
Vẽ đồ thị của hàm dưới dấu tích phân và tích phân bất định của nó trên cùng một màn hình rồi kiểm tra tính hợp lý của kết quả.

\bt{43} Tìm thể tích khối thu được bằng cách quay quanh trục $x$ miền giới hạn bởi các đường cong $y=9/(x^2+9)$, $y=0$, $x=0$ và $x=3$.

\bt{44} Tìm thể tích khối thu được bằng cách quay quanh đường thẳng $x=1$ miền nằm dưới đường cong $y=x\sqrt{1-x^2}$, $0\le x\le1$.

\par\noindent\begin{minipage}[t]{\linewidth}
\bt{45} (a) Dùng phép thế lượng giác để kiểm chứng rằng
\[\int_0^x\sqrt{a^2-t^2}\,dt=\frac12a^2\sin^{-1}(x/a)+\frac12x\sqrt{a^2-x^2}.\]
(b) Dùng hình vẽ để đưa ra cách diễn giải lượng giác cho cả hai số hạng ở vế phải của phương trình trong phần (a).
\fig[.29]{exercise45}{Hình cho Bài tập 45.}

\end{minipage}\par
\bt{46} Parabol $y=\frac12x^2$ chia hình tròn $x^2+y^2\le8$ thành hai phần. Tìm diện tích của mỗi phần.

\bt{47} Một khối xuyến được tạo thành bằng cách quay đường tròn $x^2+(y-R)^2=r^2$ quanh trục $x$. Tìm thể tích bên trong khối xuyến.

\par\noindent\begin{minipage}[t]{\linewidth}
\bt{48} Một thanh tích điện dài $L$ tạo ra tại điểm $P(a,b)$ một điện trường được cho bởi
\[E(P)=\int_{-a}^{L-a}\frac{\lambda b}{4\pi\varepsilon_0(x^2+b^2)^{3/2}}\,dx,\]
trong đó $\lambda$ là mật độ điện tích trên một đơn vị chiều dài của thanh và $\varepsilon_0$ là điện thẩm của chân không (xem hình). Tính tích phân để xác định một biểu thức cho điện trường $E(P)$.
\fig[.36]{exercise48}{Hình cho Bài tập 48.}

\end{minipage}\par
\par\noindent\begin{minipage}[t]{\linewidth}
\bt{49} Tìm diện tích miền hình lưỡi liềm (gọi là \textit{lune}) giới hạn bởi các cung tròn có bán kính $r$ và $R$. (Xem hình.)
\fig[.24]{exercise49}{Hình cho Bài tập 49.}

\end{minipage}\par
\bt{50} Một bể chứa nước có dạng hình trụ với đường kính $10$ ft. Bể được đặt sao cho các thiết diện tròn thẳng đứng. Nếu độ sâu của nước là $7$ ft thì bao nhiêu phần trăm tổng dung tích đang được sử dụng?
\end{document}
